\chapter{Conclusions and Discussion}
\label{chapter:conclusion_discussion}

\section{Synthesis of Findings}

This study applied Cooke's Classical Model to combine seven experts' uncertainty assessments into performance-weighted forecasts for the Dutch housing market in 2027. The application is appropriate because the outcome depends on interacting demographic, regulatory, financial, environmental, and labour constraints that historical extrapolation alone cannot represent. Calibration variables provided empirical control: experts were not weighted by credentials or confidence, but by statistical accuracy and informativeness on quantities with known realizations.

The optimized Decision Maker places the median number of completed dwellings at \DwellingsMedian{}, with a 90\% interval of \DwellingsFive{}--\DwellingsNinetyFive{}. Because even the upper bound is below 100,000, the panel assigns little probability to the government's annual construction ambition being achieved in 2027. The median forecasts of \PermitsMedian{} permits, an Amsterdam transaction price of \euro\PriceMedian, and \RentMedian\% free-sector rent growth point in the same direction: supply remains constrained while affordability pressure persists.

The comparison with equal weighting strengthens the interpretation. Central forecasts are similar under both schemes, but equal weighting creates a much higher upper tail for permits. The optimized DM achieves the higher combined calibration-information score and limits the influence of two strongly miscalibrated respondents. The substantive conclusion is therefore not an artefact of a single expert, although the precise tail widths remain sensitive to the aggregation rule.

\section{Critical Reflection and Methodological Limitations}

\subsection*{Panel composition and sample size}
The panel is small and was recruited mainly through the researcher's academic and professional network. It contains valuable practical housing expertise, but it is not representative of every relevant institution. Developers, mortgage lenders, municipalities, housing associations, and national policymakers may hold systematically different information. With only seven experts, individual scores also have a visible effect on the pooled distribution. Future work should recruit a larger, deliberately stratified panel.

\subsection*{Ten seed variables and finite-sample calibration}
The elicitation used ten seed variables to control participant burden. This reduces fatigue but limits the resolution of the calibration test: a single realization changes a bin frequency by ten percentage points. The chi-square approximation is also less reliable with such a small number of seeds and expected tail counts of only 0.5. The reported calibration scores should therefore be treated as useful performance indicators rather than precise population parameters. More seeds, or out-of-sample validation as additional realizations become available, would provide a stronger test.

\subsection*{Background measures and dependence}
Information scores depend on the intrinsic ranges and the choice between uniform and log-uniform backgrounds. The sensitivity analysis shows that the selected threshold, dominant experts, and dwelling median are stable under the tested alternatives, although this cannot cover every defensible specification.

The Classical Model scoring calculation also treats seed realizations as independent. Housing indicators such as mortgage debt, transaction prices, transactions, and homelessness share macroeconomic drivers, so ten questions do not necessarily provide ten independent pieces of calibration evidence. Positive dependence can make the panel appear to have more independent calibration evidence than it truly has and can overstate the effective precision of the scores. Future work could elicit or estimate a joint dependence structure after establishing the marginal distributions. Copulas, including vine constructions assembled from flexible pair-copulas, can represent dependence separately from the univariate margins; structured protocols for eliciting dependence are reviewed by Werner et al.~\cite{wernerEtAl2017}. With only seven experts and ten seeds, estimating such a structure here would add more parameters than the data can credibly support, so independence remains an explicit simplifying assumption rather than an unnoticed one.

\subsection*{Remote elicitation and data quality}
The survey was completed remotely. A bespoke interface and training examples reduced confusion, but the setting did not prevent respondents from consulting external sources. Two response patterns also reveal possible unit interpretation problems and substantial overconfidence. Rather than correcting these answers after the fact, the analysis retained the elicited values and allowed calibration to penalize them. A future elicitation should display units beside every input, validate plausible scales, conduct a short comprehension check, and record any clarifications before submission.

\subsection*{Scope of the forecast}
The four target distributions express expert uncertainty, not causal estimates. They do not identify how much of the forecast is attributable to nitrogen regulation, interest rates, labour shortages, or permitting delays. Nor can they anticipate structural breaks after elicitation. The results should inform scenario planning and policy discussion alongside econometric forecasts, not replace them.

\section{Final Assessment}

The study's strongest contribution is methodological transparency: known seed realizations discipline the aggregation, alternative Decision Makers are compared on the same performance criteria, and the final statements retain their probabilistic meaning. The analysis supports a clear but appropriately qualified conclusion: the expert panel expects the Dutch housing shortage and affordability pressure to persist through 2027, and regards the 100,000-dwelling ambition as highly unlikely under current conditions.
